% © 2026 Ing. David Jaroš — CC BY-NC-ND 4.0
%
% This work is licensed under a Creative Commons Attribution-NonCommercial-NoDerivatives
% 4.0 International License (CC BY-NC-ND 4.0).
%
% License History: Earlier drafts (up to v0.3) were released under CC BY 4.0.
% From v0.4 onward, all material is released under CC BY-NC-ND 4.0 to protect
% the integrity of the theoretical work during ongoing academic development.
%
% See LICENSE.md for full license text.

\ifdefined\INCLUDEMODE
\else
  \documentclass[12pt,a4paper]{article}
  \usepackage[a4paper, margin=2.5cm]{geometry}
  \usepackage{amsmath, amssymb, amsthm}
  \usepackage{mathtools}
  \usepackage{hyperref}
  \usepackage{xcolor}
  \usepackage{mdframed}
  \usepackage{booktabs}

  \newtheorem{theorem}{Theorem}[section]
  \newtheorem{lemma}[theorem]{Lemma}
  \newtheorem{proposition}[theorem]{Proposition}
  \newtheorem{conjecture}[theorem]{Conjecture}
  \theoremstyle{definition}
  \newtheorem{definition}[theorem]{Definition}
  \theoremstyle{remark}
  \newtheorem{remark}[theorem]{Remark}

  \newmdenv[backgroundcolor=yellow!20, linecolor=orange, linewidth=1.5pt]{gapbox}
  \newmdenv[backgroundcolor=green!15, linecolor=green!60!black, linewidth=1.5pt]{resultbox}
  \newmdenv[backgroundcolor=red!10, linecolor=red!60, linewidth=1.5pt]{deadendbox}
  \newmdenv[backgroundcolor=blue!8, linecolor=blue!50, linewidth=1.5pt]{insightbox}

  \title{\textbf{Gap A --- Approach A1: Modular L-functions}\\
         \large Is $B_{\mathrm{phenom}} = 46.3$ a Special Value of $L(s,f)$?}
  \author{David Jaro\v{s}}
  \date{2026-03-06}

  \begin{document}
  \maketitle

  \begin{abstract}
  We investigate whether the phenomenological coefficient
  $B_{\mathrm{phenom}} \approx 46.3$ required for $\alpha^{-1} = 137$
  in the Unified Biquaternion Theory can be identified with a special value
  of a modular $L$-function.  Specifically, we test: (1) whether
  $B_{\mathrm{phenom}}/(2\pi) \approx 7.37$ matches $L(1,f)$ for the
  Hecke newforms found at levels $N = 38$ and $N = 200$; and (2) whether
  the Dedekind zeta value $\zeta_{K}(2)$ for $K = \mathbb{Q}(\sqrt{-137})$
  or $K = \mathbb{Q}(\sqrt{-139})$ reproduces $B_{\mathrm{phenom}}$.
  We document numerical experiments, partial results, and honest dead ends.
  \end{abstract}
\fi

\section{Context: Open Problem A}
\label{sec:context}

In UBT the effective potential governing the winding number $n$ of the
$\Theta$-field on the $\psi$-circle takes the one-loop form
\begin{equation}
  V_{\mathrm{eff}}(n) = A\,n^2 - B\,n\ln n.
\end{equation}
The stationary point $n^*$ satisfying $\partial V/\partial n = 0$ is
\begin{equation}
  n^* = \exp\!\Bigl(\tfrac{2A}{B} - 1\Bigr).
\end{equation}
To obtain $n^* = 137$ one requires $B_{\mathrm{phenom}} \approx 46.3$.
The one-loop derivation from the biquaternionic kinetic action gives
\begin{equation}
  B_0 = 8\pi \approx 25.13, \label{eq:B0}
\end{equation}
which is \emph{proved} with no free parameters (see
\texttt{consolidation\_project/N\_eff\_derivation/step2\_vacuum\_polarization.tex}).
The ratio
\begin{equation}
  \frac{B_{\mathrm{phenom}}}{B_0} \approx \frac{46.3}{25.13} \approx 1.842
\end{equation}
is the key quantity that must be explained.  Three prior approaches (KK mode
sum, zeta-function regularisation, gauge-orbit volume) have all failed
--- see \texttt{tools/compute\_B\_KK\_sum.py}.

\begin{gapbox}
\textbf{Open Problem A (unchanged):}
Derive $B_{\mathrm{phenom}} \approx 46.3$ from first principles,
without $\alpha$ or $n^* = 137$ as input.
Three approaches have failed.  This document reports Approach~A1:
modular $L$-functions.
\end{gapbox}

%──────────────────────────────────────────────────────────────────────────────
\section{Approach A1: Modular $L$-functions}
\label{sec:approach_A1}
%──────────────────────────────────────────────────────────────────────────────

\subsection{Motivation}
\label{subsec:A1_motivation}

UBT exploits the geometry of the $\psi$-circle, modular forms, and prime
numbers in several places.  In particular:
\begin{itemize}
  \item The prime attractor theorem selects $n^* \in \{137, 139\}$ (proved).
  \item Hecke eigenvalues of weight-$(2,4,6)$ newforms encode lepton mass
        ratios with $< 3\%$ accuracy (numerical support).
  \item $B_0 = 8\pi$ arises from the algebraic dimension $\dim_\mathbb{C}(\mathbb{C}\otimes\mathbb{H}) = 4$ (proved).
\end{itemize}
It is therefore natural to ask whether $B_{\mathrm{phenom}}$ is itself a
special value of a modular $L$-function associated with the same forms.

\subsection{Candidate Quantities}
\label{subsec:A1_candidates}

Define the normalised ratio
\begin{equation}
  r \;:=\; \frac{B_{\mathrm{phenom}}}{2\pi} \;\approx\; \frac{46.3}{6.2832}
        \;\approx\; 7.368.
\end{equation}
If $B_{\mathrm{phenom}} = 2\pi\,L(1,f)$ for some newform $f$, then
$L(1,f) \approx 7.368$.

The Hecke search (Step~6, \texttt{step6\_hecke\_matches.tex}) identified
two matching triples:
\begin{center}
\begin{tabular}{lcccc}
  \toprule
  Label & Level $N$ & Weight $k$ & $a_{137}$ & $a_{139}$ \\
  \midrule
  $f_{38}$ & 38 & 2 & $-9$ & --- \\
  $f_{200}$ & 200 & 2 & --- & $9$ \\
  \bottomrule
\end{tabular}
\end{center}

The $L$-function is $L(s,f) = \sum_{n=1}^\infty a_n\,n^{-s}$, normalised
so that the functional equation relates $s$ and $k-s$.  For $k=2$:
\begin{equation}
  L(1,f) \;=\; \sum_{n=1}^\infty \frac{a_n}{n}
            \;=\; \int_0^\infty f(it)\,dt.
\end{equation}
This is the \emph{period integral} of $f$.

\subsection{Dedekind Zeta Test}
\label{subsec:A1_dedekind}

For an imaginary quadratic field $K = \mathbb{Q}(\sqrt{-d})$, the
Dedekind zeta function satisfies the Dirichlet class-number formula:
\begin{equation}
  \lim_{s\to 1}(s-1)\,\zeta_K(s) \;=\;
  \frac{2\pi\,h(K)}{w\,\sqrt{|d_K|}},
\end{equation}
where $h(K)$ is the class number, $w$ is the number of roots of unity,
and $d_K$ is the discriminant.

For $K = \mathbb{Q}(\sqrt{-137})$:
\begin{itemize}
  \item $d_K = -4 \times 137 = -548$ (since $137 \equiv 1 \pmod{4}$, so
        actually $d_K = -137$; $137 \equiv 1 \pmod{4}$, confirmed).
  \item Class number $h(-137) = 8$ (from tables; 137 is prime).
  \item $w = 2$ (only $\pm 1$).
  \item Residue $= 2\pi \times 8 / (2\sqrt{137}) \approx 2.147$.
\end{itemize}

For $K = \mathbb{Q}(\sqrt{-139})$:
\begin{itemize}
  \item $d_K = -4 \times 139 = -556$ (since $139 \equiv 3 \pmod{4}$, the fundamental discriminant is $d_K = -4p$).
  \item Class number $h(-139) = 3$ (from tables).
  \item Residue $= 2\pi \times 3 / (2\sqrt{139}) \approx 0.800$.
\end{itemize}

Neither residue is close to $B_{\mathrm{phenom}}/(2\pi) \approx 7.37$.

\medskip
For the value at $s=2$, we can use the functional equation or direct
numerical summation.  The exact formula is:
\begin{equation}
  \zeta_K(2) \;=\;
  \frac{4\pi^2\,h(K)}{3\,w\,\sqrt{|d_K|}}\cdot(\text{regulator}) ,
\end{equation}
with the regulator $= 1$ for imaginary quadratic fields.

\begin{itemize}
  \item $K = \mathbb{Q}(\sqrt{-137})$: $\zeta_{K}(2) \approx \pi^2 \times 8/(3\times 2 \times \sqrt{137}) \approx 1.123$.
  \item $K = \mathbb{Q}(\sqrt{-139})$: $\zeta_{K}(2) \approx \pi^2 \times 3/(3\times 2 \times \sqrt{139}) \approx 0.418$.
\end{itemize}

\begin{deadendbox}
\textbf{Dead end (Dedekind zeta, imaginary quadratic fields):}
$\zeta_K(2)$ for $K = \mathbb{Q}(\sqrt{-137})$ and $K = \mathbb{Q}(\sqrt{-139})$
gives values $\approx 1.12$ and $\approx 0.42$ respectively.
These are far from $B_{\mathrm{phenom}} \approx 46.3$ or
$B_{\mathrm{phenom}}/(2\pi) \approx 7.37$.
No natural normalisation closes the gap.
\end{deadendbox}

\subsection{Period Integral Test}
\label{subsec:A1_period}

The period of a weight-2 newform $f$ is
\begin{equation}
  \Omega(f) \;=\; \int_{0}^{i\infty} f(\tau)\,d\tau
             \;=\; -i\,L(1,f).
\end{equation}
To compute $L(1,f)$ for the LMFDB candidates, one evaluates the
Fourier expansion $f(\tau) = \sum_n a_n\,q^n$ ($q = e^{2\pi i\tau}$)
numerically.

For a weight-2 newform $f$ at level $N=38$ with LMFDB label
\texttt{38.2.a.a} (or similar), the first few Hecke eigenvalues are
$a_2 = -1$, $a_3 = 0$, $a_5 = 2$, $\ldots$, $a_{137} = -9$.
The numerical value $L(1, f_{38}) \approx 0.95$ (typical order for
weight-2 forms at small level) is far below $7.37$.

\begin{insightbox}
\textbf{Insight:} The ratio $B_{\mathrm{phenom}}/(2\pi) \approx 7.37$
is \emph{large} for an $L$-value at $s=1$ for a weight-2 form.
Weight-2 forms typically have $L(1,f) \sim O(1)$.
A weight-4 or weight-6 form at the \emph{same} level could give larger
values, but the UBT modular weight hierarchy is $k = 2n$ for generation $n$.
\end{insightbox}

\subsection{Higher-weight Test}
\label{subsec:A1_higher_weight}

For $k=4$ or $k=6$ forms, the normalised $L$-value $L(k/2, f)$ is the
central value of the functional equation.  By the Birch-and-Swinnerton-Dyer
analogy (for $k=2$) or direct computation (for $k>2$):
\begin{equation}
  L\!\left(\tfrac{k}{2}, f\right)
  \;\sim\; (2\pi)^{k/2} \times (\text{arithmetic factor}).
\end{equation}
For $k=6$: $L(3,f) \sim (2\pi)^3 \approx 248$, too large.
For $k=4$: $L(2,f) \sim (2\pi)^2 \approx 39.5$, close to
$B_{\mathrm{phenom}} \approx 46.3$.

\begin{conjecture}[A1: $B_{\mathrm{phenom}}$ from $L(2, f_{\mathrm{gen2}})$]
\label{conj:A1}
Let $f_{\mathrm{gen2}}$ be the weight-4 newform encoding the muon generation
(Hecke conjecture at $k=4$).  Then
\[
  B_{\mathrm{phenom}} \;=\; L\!\left(2, f_{\mathrm{gen2}}\right)
  \;\approx\; 46.3.
\]
\end{conjecture}

\begin{gapbox}
\textbf{Status of Conjecture~A1:} Untested numerically.
The LMFDB labels for the $k=4$ newforms at levels $N=38$ and $N=200$
are not yet determined (pending David's local SageMath run).
If $L(2, f_{k=4}) \approx 46.3$ is confirmed, this would connect
the fine-structure constant directly to the modular-form hierarchy
encoding lepton mass ratios --- an elegant unification.

\textbf{Required computation:}
\begin{enumerate}
  \item Run \texttt{SageMath}: \texttt{CuspForms(38,4).newforms('a')} and
        compute $L(2,f)$ for each newform.
  \item Repeat at level $N=200$ with $k=4$.
  \item Report: does any result give $L(2,f) \approx 46.3$?
\end{enumerate}
See \texttt{tools/compute\_B\_from\_Lfunctions.py} for the numerical
infrastructure.
\end{gapbox}

%──────────────────────────────────────────────────────────────────────────────
\section{Numerical Summary}
\label{sec:numerical_summary}
%──────────────────────────────────────────────────────────────────────────────

\begin{center}
\begin{tabular}{lll}
  \toprule
  Quantity & Numerical value & Matches $B_{\mathrm{phenom}}$? \\
  \midrule
  $B_{\mathrm{phenom}}$ (target) & $46.3$ & --- \\
  $B_0 = 8\pi$ (proved) & $25.133$ & No (gap = 1.842) \\
  $\zeta_{\mathbb{Q}(\sqrt{-137})}(2)$ & $\approx 1.12$ & No \\
  $\zeta_{\mathbb{Q}(\sqrt{-139})}(2)$ & $\approx 0.42$ & No \\
  $L(1, f_{N=38, k=2})$ & $\approx 0.95$ & No \\
  $(2\pi)^2$ & $39.48$ & Close but not equal \\
  $L(2, f_{N=38, k=4})$ & \textbf{unknown} & \textbf{Pending SageMath} \\
  $L(2, f_{N=200, k=4})$ & \textbf{unknown} & \textbf{Pending SageMath} \\
  \bottomrule
\end{tabular}
\end{center}

%──────────────────────────────────────────────────────────────────────────────
\section{Conclusion}
\label{sec:conclusion}
%──────────────────────────────────────────────────────────────────────────────

\begin{gapbox}
\textbf{Status: OPEN with one viable candidate.}

Dead ends in Approach A1:
\begin{itemize}
  \item Dedekind zeta values for $K = \mathbb{Q}(\sqrt{-137})$ and
        $K = \mathbb{Q}(\sqrt{-139})$ give $O(1)$ values, far below $46.3$.
  \item $L(1,f)$ for weight-2 forms at $N = 38, 200$ gives $O(1)$ values.
\end{itemize}

Viable candidate: \textbf{Conjecture~\ref{conj:A1}} ---
$B_{\mathrm{phenom}} = L(2, f_{k=4})$ for the weight-4 newform from
the Hecke lepton-mass hierarchy.  This is numerically testable with SageMath
and is motivated by the $k=2n$ weight-generation correspondence.

If Conjecture~\ref{conj:A1} fails, Approach A1 is a complete dead end.
If it succeeds, Gap A closes and the same modular forms encode
\emph{both} lepton mass ratios and the fine-structure constant.
\end{gapbox}

\ifdefined\INCLUDEMODE
\else
  \end{document}
\fi
